\begin{abstract}
Large language models (LLMs) increasingly replace human annotators in computational social science, yet their error rates can vary by up to $5\times$ across demographic and text-feature subgroups.
When these noisy labels serve as effect modifiers in heterogeneous treatment effect (HTE) analysis, the standard remedy, estimating a single confusion matrix and applying its inverse uniformly, is widely trusted to remove bias.
We prove it can do the opposite.
We derive an exact bias identity showing that global correction introduces a subgroup-level bias proportional to the interaction of error-rate deviation and effect heterogeneity, a product structure invisible to any pooled analysis.
This identity reveals that correction reverses estimated causal effect signs in 14.9\% of configurations and amplifies bias across all nine real LLMs tested.
A diagnostic corollary flags harmful correction from error rates alone, without estimating treatment effects.
Guided by the identity, we propose Error-Corrected HTE (EC-HTE), a per-subgroup correction via hierarchical Bayesian confusion matrix estimation that reduces bias by $7.4\times$ under extreme heterogeneity, with clear operating conditions (Section~\ref{sec:diagnostic}), validated across 11 LLM--dataset combinations spanning four NLP tasks\footnote{Code and data: \url{https://anonymous.4open.science/r/ec_hte_llm_annotation-4CC8/}.}.
\end{abstract}